\documentclass[12pt]{article}
\usepackage[utf8]{inputenc}
\usepackage[T1]{fontenc}
\usepackage{hyperref}
\title{Taveta Language Resources for K-12 Curriculum Development}
\author{Compiled by Claude Sonnet 5 (claude-sonnet-5) for Andr\'as Kornai}
\date{}
\begin{document}
\maketitle

\section{Introduction}
Taveta (kiTaveta, ISO 639-3 tvs) is a Northeast Coast Bantu language of Kenya. It is closely related to Pare. Speaker counts in the sources differ: Wikipedia cites about 21,000 first-language speakers (2009 census), and Joshua Project lists about 37,000 across Kenya and Tanzania. Little material exists beyond scripture and general reference pages. No school textbooks, primers or dictionaries turned up in the searches.

\section{Primary Educational Resources}
None were found in search results.

\section{Linguistic and Reference Materials}
\begin{itemize}
\item Wikipedia, Taveta language: \url{https://en.wikipedia.org/wiki/Taveta_language}
\end{itemize}
No dictionary or grammar turned up in the searches.

\section{Religious and Community Resources}
\begin{itemize}
\item Taveta Bible text on YouVersion: \url{https://www.bible.com/versions/3353}
\item Joshua Project language profile (reports no complete Bible and no audio Bible): \url{https://www.joshuaproject.net/index.php/languages/tvs}
\end{itemize}
Anglican missionary translations date from the 1890s, and a New Testament was completed in 1906 according to the search summary.

\section{Digital and Audio Resources}
No audio Bible is reported by Joshua Project. No other digital resources were found.

\section{Curriculum Development Context}
The language is not known to be taught in schools. A curriculum would have to start from scratch, using the old scripture translations as a text base and new primers developed with the community. Any orthography would need to be checked against current usage.
\end{document}